%%%%%%%%%%%%%%%%%%%%%%%%%%%%%%%%%%%%%%%%%%%%%%%%%%%%%%%%%%%%%%%%%%%%%%%%
% Plantilla TFG/TFM
% Universidad de Murcia. Facultad de Informática
% Realizado por: José Manuel Requena Plens
% Modificado: Pablo José Rocamora Zamora
% Contacto: pablojoserocamora@gmail.com
%%%%%%%%%%%%%%%%%%%%%%%%%%%%%%%%%%%%%%%%%%%%%%%%%%%%%%%%%%%%%%%%%%%%%%%%


\chapter{Introducción}

\section{Contexto y Motivación}

La democratización de los Grandes Modelos de Lenguaje (LLMs) ha marcado un cambio de paradigma en la generación y recuperación de información. Sin embargo, esta misma capacidad para generar texto a gran velocidad se convierte en un problema cuando el entorno informativo está lleno de desinformación. La inteligencia artificial generativa desempeña un papel dual en el panorama mediático actual: facilita la creación rápida y la diseminación dirigida de contenido sintético pero también presenta nuevas oportunidades tecnológicas para la detección, la educación mediática y la verificación automatizada \cite{mencia2025mapping}.

Ante esta realidad, la validación manual de hechos resulta cada vez más insuficiente debido a la velocidad sin precedentes con la que la información y la desinformación se propagan \cite{guo2022survey}. Por ello, resulta imperativa la transición hacia herramientas de \textit{Fact-Checking} automatizado que se apoyen en el procesamiento del lenguaje natural y la inteligencia artificial para predecir la veracidad de las afirmaciones y asistir a los profesionales de la información frente a volúmenes masivos de datos \cite{guo2022survey}.

\section{Planteamiento del Problema: Las limitaciones de los LLMs}

Si bien los LLMs han revolucionado el procesamiento del lenguaje, depender exclusivamente de su memoria interna para tareas de verificación de hechos plantea riesgos inasumibles. Una de las limitaciones fundamentales es la degradación temporal de su conocimiento. Los modelos asimilan el conocimiento del mundo en sus parámetros durante la fase de preentrenamiento, pero esta información puede quedar rápidamente desactualizada a medida que la realidad cambia \cite{jang2022towards}. Actualizar constantemente estos modelos es un desafío técnico, ya que el reentrenamiento continuo para adquirir nueva información suele provocar un ``olvido catastrófico" del conocimiento previamente adquirido, o requerir expansiones de parámetros computacionalmente ineficientes \cite{jang2022towards}. En consecuencia, su base de conocimiento interno tiende a permanecer estática tras el entrenamiento inicial.

A esta falta de actualización se suma el problema más crítico en el ámbito de la veracidad: el fenómeno de la alucinación. En el contexto de los LLMs, la alucinación se define como la generación de contenido que resulta fluido y sintácticamente correcto, pero que es factualmente inexacto o no está respaldado por ninguna evidencia externa \cite{alansari2025hallucination}. Esta tendencia intrínseca a inventar datos plausibles perjudica directamente la fiabilidad y la confianza en estos sistemas, volviéndolos inadecuados para operar de forma aislada en dominios que exigen una precisión factual absoluta, como es el caso del \textit{Fact-Checking} \cite{alansari2025hallucination}.

\section{Propuesta de Solución: Generación Aumentada por Recuperación (RAG)}

Para mitigar las alucinaciones de los LLMs, este trabajo propone la implementación de una arquitectura de Generación Aumentada por Recuperación (\textit{Retrieval-Augmented Generation}, RAG). Conceptualmente, RAG combina dos tipos de memoria \cite{lewis2020rag}: una paramétrica, formada por el propio modelo generativo preentrenado que almacena conocimiento en sus pesos, y una no paramétrica, formada por un índice externo de documentos consultable mediante un recuperador, cuyo contenido puede actualizarse sin necesidad de reentrenar el modelo.

La arquitectura RAG actúa como un puente tecnológico que vincula la capacidad de razonamiento y la fluidez narrativa de los LLMs con una ``fuente de verdad" externa. En el contexto de este proyecto, dicha fuente está representada por un corpus masivo de noticias de diversos medios y temas. Al realizar una consulta, el sistema recupera las evidencias más relevantes antes de la generación de la respuesta, lo que permite anclar la salida del modelo en hechos contrastables y reducir drásticamente la probabilidad de alucinación factual.

\section{Alcance y Enfoque del Proyecto}

Este Trabajo Fin de Grado aborda la construcción y validación de un sistema de verificación de información escalable y reproducible. Para ello, se parte de un corpus de aproximadamente 800.000 noticias indexadas y gestionadas mediante el motor de búsqueda Elasticsearch.

El núcleo del trabajo se centra en una comparativa multidimensional. Por un lado, se comparan dos estrategias de recuperación: la búsqueda léxica clásica con BM25 frente a la recuperación semántica mediante el modelo de \textit{embeddings} BGE-M3 \cite{chen2024bge}. Por otro lado, se analiza el impacto del tamaño del modelo generador, comparando el rendimiento de modelos de parámetros masivos, como Llama 3.3 70B entre ellos, frente a modelos de lenguaje pequeños (SLM) optimizados para su ejecución local, como es el caso de Phi-3 \cite{abdin2024phi3}.

Finalmente, dada la complejidad de corregir manualmente el elevado volumen de respuestas generadas, se implementa el paradigma \textit{LLM-as-a-judge} \cite{zheng2023judging}, utilizando un modelo \texttt{gpt-oss-120b} \cite{openai2025gptoss} como juez automático que simula la corrección humana y facilita la identificación de errores en las respuestas generadas por las distintas combinaciones de recuperador y generador.

\section{Estructura de la Memoria}

La memoria se organiza de manera secuencial en cinco capítulos que reflejan el desarrollo del proyecto:

\begin{itemize}
    \item \textbf{Capítulo 2: Estado del arte.} Se exploran los fundamentos teóricos del TFG. Se detallan los mecanismos de funcionamiento de los \textit{Transformers}, los modelos de \textit{embeddings} y el funcionamiento de los motores de búsqueda tanto léxicos como vectoriales, estableciendo la base técnica necesaria para comprender la arquitectura RAG. Además, se analizan las capacidades de los SLMs y se explica el paradigma \textit{LLM-as-a-judge} como método de evaluación automatizada.
    
    \item \textbf{Capítulo 3: Análisis de objetivos y metodología.} Se definen formalmente los objetivos generales y específicos que persigue el proyecto y se introduce el desarrollo en siete fases del mismo, abarcando desde la ingesta inicial del corpus hasta la síntesis final del sistema óptimo.
    
    \item \textbf{Capítulo 4: Diseño y resolución del trabajo realizado.} Es el capítulo más extenso de la memoria. En él se detalla la infraestructura utilizada en el desarrollo, el proceso de ingeniería de datos (ETL) sobre el corpus de noticias y el diseño de cada uno de los sistemas a evaluar. Se incluye el análisis de resultados obtenidos tras evaluar el sistema con 220 preguntas, comparando el rendimiento de las distintas combinaciones de recuperadores y generadores de diferentes maneras. Se razonan los hallazgos obtenidos y en base a eso se construye un sistema final que también es evaluado.
    
    \item \textbf{Capítulo 5: Conclusiones y trabajos futuros.} Se sintetizan los hallazgos principales del TFG, validando el cumplimiento de los objetivos iniciales y exponiendo las limitaciones del proyecto. Finalmente, se reflexiona sobre el impacto ético del uso de IA en el periodismo y se proponen líneas futuras para seguir optimizando la precisión y escalabilidad de los sistemas de verificación.
\end{itemize}